\documentclass[a4paper,12pt]{article}
\usepackage{pdfpages}
\usepackage[cp1252]{inputenc}
\usepackage{lipsum}
\usepackage{listings}
\usepackage{pdflscape}
\usepackage{graphicx}
\usepackage[margin=1.5cm]{geometry} % Verkleinert die R�nder
\definecolor{okgreen}{rgb}{0.0, 0.5, 0.0}

\lstdefinestyle{cppstyle}{
  language=C++,
  basicstyle=\ttfamily\tiny,
  keywordstyle=\color{blue},
  commentstyle=\textcolor{okgreen},
  stringstyle=\color{red},
  numbers=left,
  numberstyle=\tiny\color{gray},
  stepnumber=1,
  breaklines=true,
  frame=single
}

\definecolor{failred}{rgb}{0.7, 0.0, 0.0}

\lstdefinelanguage{TestOutput}{
    morekeywords={},
    morecomment=[l]{//},
    morestring=[b]",
    sensitive=false,
}

\lstdefinestyle{teststyle}{
    language=TestOutput,
    basicstyle=\ttfamily\small,
    keywordstyle=\color{black},
    commentstyle=\color{gray},
    stringstyle=\color{black},
    showstringspaces=false,
    breaklines=true,
    frame=single,
    numbers=left,
    numberstyle=\tiny\color{gray},
    postbreak=\mbox{\textcolor{red}{$\hookrightarrow$}\space},
    literate={OK}{{\textcolor{okgreen}{OK}}}2
             {Fail}{{\textcolor{failred}{Fail}}}4,
}

\begin{document}
\includepdf[pages=1-3] {Angabe_Uebung8.pdf}
\newpage
\section{Beispiel 1 Shift Template}
\subsection{L�sungsidee}
Die vorliegende Funktion bietet eine generische L�sung zum Linksschieben von Elementen in beliebigen Containern mittels Iteratoren. Dabei werden die ersten count Elemente entfernt, die restlichen nach vorne verschoben und die freigewordenen Positionen am Ende mit einem festen Wert aufgef�llt. Durch die Verwendung von Templates ist die Funktion flexibel und f�r verschiedene Containertypen wie std::vector oder std::list einsetzbar, sofern deren Iteratoren grundlegende Operationen wie Inkrementieren und Dereferenzieren unterst�tzen. Die Implementierung erfolgt in-place, wodurch keine zus�tzlichen Speicherressourcen ben�tigt werden und die Effizienz gewahrt bleibt.
\subsection{ShiftTemplate.hpp}
\lstinputlisting[style=cppstyle]{Uebung8/Uebung8/ShiftTemplate.hpp}
\newpage
\section{Beispiel 2 Count Pairs}
\subsection{L�sungsidee}
Die Funktion \texttt{CountPairs} bietet eine generische M�glichkeit, echte Paare in einem Container zu z�hlen � also exakt zwei unmittelbar aufeinanderfolgende, gleiche Elemente. Durch die Nutzung von \texttt{std::adjacent\_find} wird effizient nach solchen Paaren gesucht, wobei l�ngere Sequenzen gleicher Werte bewusst ausgeschlossen werden. Die Funktion arbeitet iteratorbasiert und ist dadurch f�r verschiedenste Containertypen flexibel einsetzbar.
\subsection{CountPairs.hpp}
\lstinputlisting[style=cppstyle]{Uebung8/Uebung8/CountPairs.hpp}
\newpage
\section{Beispiel 3 Set Operations}
\subsection{L�sungsidee}
Die bereitgestellten Template-Funktionen \texttt{Union}, \texttt{Intersection} und \texttt{Difference} erm�glichen grundlegende Mengenoperationen auf beliebigen Containern mittels Iteratoren. Ziel ist es, die Vereinigungs-, Schnitt- und Differenzmenge zweier Container generisch zu berechnen und das Ergebnis �ber einen Zieliterator (\texttt{begin3}) auszugeben. Dabei wird f�r jede Operation sichergestellt, dass doppelte Elemente vermieden werden, indem diese mithilfe von \texttt{std::find} �berpr�ft werden. Die Funktionen arbeiten vollst�ndig iteratorbasiert und lassen sich somit flexibel f�r verschiedene Containertypen wie \texttt{std::vector}, \texttt{std::list} oder \texttt{std::deque} einsetzen. 
\subsection{SetOperations.hpp}
\lstinputlisting[style=cppstyle]{Uebung8/Uebung8/SetOperations.hpp}
\newpage
\begin{landscape}
\subsection{Testtreiber main.cpp}
\lstinputlisting[style=cppstyle]{Uebung8/Uebung8/main.cpp}
\end{landscape}
\newpage
\begin{landscape}
\subsection{Testfunktionen Test.hpp}
\lstinputlisting[style=cppstyle]{Uebung8/Uebung8/Test.hpp}
\newpage
\subsection{Tests}
F�r die Tests der Templates wurde darauf geachtet, dass diese auf verschieden Container und Datentypen angewendet wurde.
Besonders der Test von verschiedenen Iteratortypen ist hier wichtig. (Forward Iterator, Random Access Iterator usw.).
Die Testf�lle sind in der Ausgabe genauer beschrieben.
\subsubsection{Testoutput:}
\lstinputlisting[style=teststyle]{Uebung8/Uebung8/TestResult.txt}
\end{landscape}
\end{document}